\documentclass[11pt,letterpaper]{article}
\usepackage[utf8]{inputenc}
\usepackage[margin=0.4in]{geometry}
\usepackage{enumitem}
\usepackage{hyperref}
\usepackage{titlesec}
\usepackage{fontawesome5}
\usepackage{xcolor}
\usepackage{needspace}
\usepackage{ragged2e}
\usepackage{tabularx}

\hypersetup{
    colorlinks=true,
    linkcolor=blue,
    filecolor=blue,
    urlcolor=blue,
    hypertexnames=false,
    pdftitle={Abhirup Mukherjee - Curriculum Vitae},
    pdfauthor={Abhirup Mukherjee},
    pdfsubject={Academic curriculum vitae, September 2026}
}
\titleformat{\section}{\normalsize\bfseries\uppercase}{}{0em}{}[\titlerule]
\titlespacing{\section}{0pt}{7pt}{3pt}
\titleformat{\subsection}{\normalsize\bfseries}{}{0em}{}
\titlespacing{\subsection}{0pt}{5pt}{2pt}
\pagenumbering{gobble}
\linespread{0.96}
\setlength{\parindent}{0pt}
\setlist[itemize]{leftmargin=*,itemsep=0pt,topsep=2pt,parsep=0pt,partopsep=0pt}
\setlength{\emergencystretch}{1em}
\raggedbottom
\clubpenalty=10000
\widowpenalty=10000
\let\cvsection\section
\renewcommand{\section}{\Needspace{7\baselineskip}\cvsection}
\newsavebox{\cvheadingbox}
\newsavebox{\cvdatebox}

\newcommand{\educationitem}[4]{%
    \Needspace{5\baselineskip}%
    \textbf{#1} \hfill \textbf{#2}\\
    #3 \hfill #4
}
\newcommand{\experienceitem}[3]{%
    \Needspace{5\baselineskip}%
    \sbox{\cvheadingbox}{\textbf{#1}, \textit{#2}}%
    \sbox{\cvdatebox}{\textbf{#3}}%
    \ifdim\dimexpr\wd\cvheadingbox+\wd\cvdatebox+1em\relax<\linewidth
        \usebox{\cvheadingbox}\hfill\usebox{\cvdatebox}%
    \else
        \textbf{#1}\hfill\usebox{\cvdatebox}\\
        \textit{#2}%
    \fi
}
\newcommand{\blackhref}[2]{\href{#1}{\color{black}#2}}
\newcommand{\skillrow}[2]{\textbf{#1} & #2\\[1pt]}

\begin{document}
\RaggedRight
\begin{center}
    {\LARGE\textbf{ABHIRUP MUKHERJEE}}\\[7pt]
    \small
    \faEnvelope\ \href{mailto:axm240026@utdallas.edu}{axm240026@utdallas.edu} \quad
    \faLinkedin\ \href{https://linkedin.com/in/mukherjeeabhirup}{mukherjeeabhirup} \quad
    \faGlobe\ \href{https://mukherjeeabhirup.github.io/}{mukherjeeabhirup.github.io} \quad
    \faGithub\ \href{https://github.com/mukherjeeabhirup}{mukherjeeabhirup}
\end{center}

\section*{Research Interests}
I'm a Ph.D. student in Computer Science at The University of Texas at Dallas (UTD), working with Prof.\ \blackhref{https://hoak.me/}{Blaine Hoak} as a Graduate Research Assistant on problems at the intersection of adversarial machine learning and computer security. Broadly, I'm interested in interpretable and robust AI models in security-critical applications.

\section*{Education}
\educationitem{The University of Texas at Dallas}{August 2026 - Present}{Doctor of Philosophy in Computer Science}{Richardson, TX}
\begin{itemize}
    \item Advisor: Prof.\ \blackhref{https://hoak.me/}{Blaine Hoak}
\end{itemize}
\vspace{2pt}
\educationitem{The University of Texas at Dallas}{August 2024 - May 2026}{Master of Science in Computer Science}{Richardson, TX}
\begin{itemize}
    \item Thesis Advisor: Dr.\ \blackhref{https://personal.utdallas.edu/~lkhan/}{Latifur Khan}
    \item Thesis: \blackhref{https://www.proquest.com/docview/3354974136}{Activation Steering for Smart Contract Vulnerability Detection}
\end{itemize}
\vspace{2pt}
\educationitem{SRM Institute of Science and Technology}{June 2019 - May 2023}{Bachelor of Technology in Electronics and Communication Engineering}{Kattankulathur, India}

\section*{Publications}
\subsection*{Conference Proceedings}
\begin{itemize}
    \item \textbf{A. Mukherjee}, L. Khan, and B. Thuraisingham, ``Activation Steering for Smart Contract Vulnerability Detection: A Calibration Mechanism Bounded by Pre-Training Knowledge,'' \textit{IEEE International Conference on Intelligence and Security Informatics (ISI)}, Cambridge, UK, 2026. Oral presentation. \href{https://mukherjeeabhirup.github.io/assets/pdf/isi_paper_poster_portrait_CCC.pdf}{[Poster]} \href{https://mukherjeeabhirup.github.io/assets/pdf/paper12-isi-2026-v3.pdf}{[Slides]}
\end{itemize}
\subsection*{Thesis}
\begin{itemize}
    \item \textbf{A. Mukherjee}, ``Activation Steering for Smart Contract Vulnerability Detection,'' \textit{Master's thesis, The University of Texas at Dallas}, 2026. \href{https://www.proquest.com/docview/3354974136}{[Thesis]}
\end{itemize}

\section*{Professional Service}
\subsection*{External Reviewer}
\begin{itemize}
    \item ACM Workshop on Artificial Intelligence and Security (AISec) - 2026
    \item USENIX Security Symposium - 2027
\end{itemize}

\section*{Research Experience}
\experienceitem{Graduate Research Assistant}{The University of Texas at Dallas, Richardson, TX}{August 2026 - Present}
\begin{itemize}
    \item Advised by Dr.\ \blackhref{https://hoak.me/}{Blaine Hoak}; working at the intersection of adversarial machine learning and computer security.
\end{itemize}
\vspace{2pt}
\experienceitem{Master's Thesis Research}{The University of Texas at Dallas}{August 2025 - May 2026}\\
\textit{Advisor: Dr.\ \blackhref{https://personal.utdallas.edu/~lkhan/}{Latifur Khan}} \quad
\href{https://mukherjeeabhirup.github.io/projects/activation-steering/}{[Project]}
\begin{itemize}
    \item Investigated activation steering for Solidity smart contract vulnerability detection, evaluating six LLMs (7B-34B parameters) on an 11-class classification task using the SmartBugs Curated dataset.
    \item Computed contrastive steering vectors from vulnerable/patched code pairs and applied them at inference time without fine-tuning.
    \item Evaluated direction-specific confidence calibration and improvements in multi-class vulnerability classification beyond few-shot prompting; this work became a first-author IEEE ISI 2026 paper.
    \item Replicated and evaluated VulnPatch using \textit{QLoRA fine-tuning} of CodeLlama on multi-GPU HPC clusters, identifying failure modes including mode collapse on small datasets and insufficient Solidity competence in the base model.
\end{itemize}
\vspace{2pt}
\experienceitem{Undergraduate Researcher}{SRM Institute of Science and Technology}{May 2022 - October 2022}\\
\textit{Advisor: Dr.\ \blackhref{https://www.linkedin.com/in/prof-dr-t-rama-rao-50197a93/}{T. Rama Rao}} \quad
\href{https://mukherjeeabhirup.github.io/projects/single-image-3d-reconstruction/}{[Project]}
\begin{itemize}
    \item Led development of a cost-effective \textit{3D reconstruction} pipeline from single 2D images for SLAM and VR applications in a three-member research team.
    \item Implemented and benchmarked feature matching algorithms (\textit{SIFT}, \textit{ORB}) and \textit{epipolar geometry} calculations using \textit{OpenCV}, \textit{PyTorch}, and \textit{Blender}.
    \item Optimized projection matrices to improve 3D reconstruction accuracy, surpassing traditional optimization baselines.
    \item Developed device-agnostic image processing techniques for rapid generation of 3D assets suitable for \textit{autonomous navigation systems}.
\end{itemize}

\section*{Work History}
\experienceitem{HPC Student Assistant}{The University of Texas at Dallas}{January 2025 - May 2026}
\begin{itemize}
    \item Accelerated scientific applications such as \textit{Siesta}, achieving \textit{10x} speedup through multi-level parallelization and batch processing using \textit{OpenMP}, \textit{MPI}, and \textit{launcher}.
    \item Resolved \textit{50+} support tickets from university researchers involving debugging and optimizing parallel scientific codes in \textit{C}, \textit{C++}, \textit{Fortran}, and \textit{Python}.
    \item Fine-tuned large language models on HPC clusters for research applications, gaining hands-on experience with \textit{distributed training} and \textit{GPU acceleration}.
    \item Assisted in identifying optimal HPC cluster configuration settings for performance boosts of up to \textit{25\%} through systematic benchmarking using custom computation kernels and industry-standard benchmarking tools.
\end{itemize}
\vspace{2pt}
\experienceitem{Product Development Engineer}{Comviva Technologies}{February 2023 - July 2024}
\begin{itemize}
    \item Developed secure payment gateway features for \textit{Qatar National Bank's (QNB)} Digital Wallet \textit{iOS application}, implementing efficient data structures for high-volume transaction processing.
    \item Architected and optimized \textit{REST APIs} across \textit{80+ new screens} using \textit{Swift} and \textit{SwiftUI}, focusing on memory-efficient data transfers for highly responsive UIs.
    \item Delivered production-ready \textit{iOS apps} at a \textit{two-week} deployment cadence, including comprehensive software testing across \textit{three} QA gates, contributing to \textit{10\%} product unit performance improvement.
    \item Implemented unit testing frameworks such as \textit{XCTest} and \textit{Quick/Nimble}, increasing code coverage by \textit{30\%}.
    \item Enhanced code quality metrics by \textit{30\%} through systematic profiling, debugging, and integration of \textit{SonarQube} static analysis in CI/CD pipelines.
\end{itemize}
\vspace{2pt}
\experienceitem{Machine Learning Research Intern}{Atos Syntel}{July 2022 - September 2022}
\begin{itemize}
    \item Optimized \textit{CNN-based} medical imaging models for mammogram classification using \textit{PyTorch} and \textit{CUDA}.
    \item Achieved \textit{91\% accuracy} on breast cancer detection while reducing inference time by \textit{40\%} through efficient memory management and batch processing.
    \item Implemented an \textit{EDA} pipeline using \textit{NumPy} and \textit{Pandas} for the Breast Cancer Wisconsin dataset.
    \item Conducted comparative analysis of deep learning (CNN) versus traditional ML approaches (Logistic Regression: \textit{97.07\%} accuracy) to identify tradeoffs between computational efficiency and accuracy in medical diagnosis applications.
    \item Engaged in discussions on bias mitigation in medical AI systems, particularly for datasets affecting minority populations.
\end{itemize}

\section*{Selected Projects}
\experienceitem{\blackhref{https://mukherjeeabhirup.github.io/projects/data-pattern-observatory/}{Data Pattern Observatory - HackUTD 2024}}{2nd Place Overall (1,100+ participants)}{November 2024}\\
\href{https://github.com/mukherjeeabhirup/data-observability-hackutd}{[Code]} \quad
\href{https://devpost.com/software/data-pattern-observatory}{[Devpost]}
\begin{itemize}
    \item Developed a data observability system for the \textit{PNC Bank} challenge to visualize patterns and identify anomalies with \textit{85\%+} accuracy as data flows across an organization.
    \item Architected a full-stack data processing pipeline using \textit{Python}, \textit{MongoDB}, \textit{Redis}, \textit{Elasticsearch}, and \textit{Docker}-containerized microservices with a \textit{Streamlit} front end.
\end{itemize}
\vspace{2pt}
\experienceitem{\blackhref{https://mukherjeeabhirup.github.io/projects/mars-lander-simulation-game/}{Mars Lander Simulation Game}}{NASA Space Apps Challenge}{October 2021}\\
\href{https://github.com/mukherjeeabhirup/TeamMarooned5}{[Code]} \quad
\href{https://mukherjeeabhirup.github.io/assets/pdf/marooned-5-slides.pdf}{[Slides]}
\begin{itemize}
    \item Led ideation and development of a Mars lander simulation game in a team of five; placed 2nd locally and was selected as a global nominee among thousands of international submissions.
\end{itemize}
\vspace{2pt}
\experienceitem{\blackhref{https://mukherjeeabhirup.github.io/projects/nodemcu-iot-smart-car/}{NodeMCU WiFi-Controlled IoT Smart Car}}{SRM Institute of Science and Technology}{July 2021}\\
\href{https://mukherjeeabhirup.github.io/assets/video/nodemcu-smart-car-web.mp4}{[Demo]}
\begin{itemize}
    \item Designed and implemented an \textit{embedded system} for a remote-controlled vehicle using an \textit{ESP8266 WiFi SoC} and \textit{L298N} motor driver.
    \item Implemented direction and speed controls in \textit{Embedded C} for a high-power motor driver (5-35V, 25W).
    \item Created local network infrastructure and an \textit{Android} application interface using \textit{Kotlin} for wireless control.
\end{itemize}

\section*{Skills}
\begin{tabularx}{\linewidth}{@{}>{\raggedright\arraybackslash}p{0.34\linewidth}@{\hspace{12pt}}>{\RaggedRight\arraybackslash}X@{}}
\skillrow{Machine Learning and Cybersecurity}{Machine Learning, Deep Learning, Computer Vision, Natural Language Processing, Model Interpretability, Activation Steering, Cybersecurity, Smart Contract Vulnerability Analysis, Secure Software Development}
\skillrow{Deep Learning Tools/Frameworks}{PyTorch, TensorFlow, Hugging Face Transformers, PyTorch Lightning, Scikit-Learn, Keras, OpenCV, Spark MLlib}
\skillrow{Programming Languages}{Python, C, C++, Swift, SwiftUI, Fortran, SQL, Bash, \LaTeX, MATLAB}
\skillrow{HPC and Parallel Computing}{NVIDIA CUDA, OpenMP, MPI, SLURM, CuPy, High Throughput Computing}
\skillrow{Big Data and Cloud}{Apache Spark, Kafka, MongoDB, Redis, Elasticsearch, Docker, AWS (EC2, S3)}
\skillrow{Development Tools}{Git/GitHub, GitLab, Docker, CI/CD (Jenkins), Postman, SonarQube, GDB, Perf, JIRA}
\textbf{Operating Systems} & Linux/Unix, Parallel Programming, POSIX Threads, Shell Scripting
\end{tabularx}

\section*{Relevant Courses}
Machine Learning $\bullet$ Deep Learning $\bullet$ Computer Vision $\bullet$ Natural Language Processing $\bullet$ Big Data Management and Analytics $\bullet$ Database Design $\bullet$ Design and Analysis of Computer Algorithms $\bullet$ Artificial Intelligence $\bullet$ Operating Systems Concepts $\bullet$ Scientific Computing with Python $\bullet$ Computer Communication Networks $\bullet$ ARM Based Embedded System Design $\bullet$ Microprocessors and Microcontrollers $\bullet$ Digital Signal Processing $\bullet$ Machine Perception with Cognition

\section*{Awards and Achievements}
\begin{itemize}
    \item \textbf{HackUTD 2024 - 2nd Place Overall} - Achieved 2nd place among 1,100+ participants for an enterprise-grade data observability system.
    \item \textbf{NASA Space Apps Challenge 2021 - 2nd Place Locally and Global Nominee} - Led ideation and development of a Mars lander simulation game in a team of five, selected as a global nominee among thousands of international submissions.
    \item \textbf{Departmental Scholarship, SRM Institute} - Awarded a merit-based scholarship in recognition of outstanding academic performance (9.46/10 CGPA) during academic year 2020-21.
\end{itemize}

\section*{Leadership Experience}
\experienceitem{Committee Head, Team Envision}{SRM Institute of Science and Technology}{July 2021 - December 2021}
\begin{itemize}
    \item Recruited and managed a team of \textit{15} members for the Electronics and Mechanical domain of the university's flagship technical festival.
    \item Maintained technical and financial documentation for operations of two technical campus clubs.
    \item Coordinated with corporate sponsors and organized technical workshops and competitions for \textit{500+} participants.
\end{itemize}

\section*{Professional Development and Certifications}
\begin{itemize}
    \item Applied Data Science with Python Specialization (IBM/Coursera)
    \item Deep Learning Applications for Computer Vision (University of Colorado Boulder/Coursera)
    \item Python for Data Science, AI and Development (IBM/Coursera)
\end{itemize}
\end{document}
